% =========================================================================
% SEQUENCE DIAGRAM FARMEASE - MODUL PETERNAKAN
% Pola: Actor -> ui:Antarmuka -> NamaKelas.method() -> db:Database
% Sintaks: \begin{call}, \begin{sdblock}, \resizebox (pgf-umlsd / TikZ)
% =========================================================================
%
% Preamble yang diperlukan:
%   \usepackage{tikz}
%   \usepackage{pgf-umlsd}
%
% =========================================================================

\documentclass{article}
\usepackage{tikz}
\usepackage{pgf-umlsd}

\begin{document}

% =========================================================================
% UC-01: MELIHAT DASBOR TERNAK
% =========================================================================

% -------------------------------------------------------------------------
% UC-01 Normal-1: Operator Ternak melihat Dasbor Ternak
% -------------------------------------------------------------------------
\item{Sequence Diagram UC-01 N-1: Operator Melihat Dasbor Ternak}

Operator Ternak membuka halaman Dasbor Ternak. Antarmuka (
\texttt{ui:DashboardView}) memanggil \texttt{GetCageStats()} pada kelas
\texttt{cs:Cage} untuk mengambil ringkasan statistik populasi ternak dan
utilitas kandang, kemudian memanggil \texttt{GetMyTasks(id, role, date)}
pada kelas \texttt{tk:Task} untuk memuat daftar tugas rutin yang dijadwalkan
pada hari tersebut. Kedua method meneruskan permintaan ke \texttt{db:Database}
dan mengembalikan data yang dibutuhkan. Setelah semua data terkumpul,
dasbor menampilkan kartu ringkasan populasi domba, daftar tugas harian, dan
widget peringatan kehamilan terdekat. Apabila Operator mengeklik salah satu
tugas pada daftar, sistem memanggil kembali \texttt{GetMyTasks()} untuk
mengambil detail tugas dan mengarahkan Operator ke formulir pencatatan yang
telah ter-\textit{prefill} secara otomatis. Alur ini disajikan pada
Gambar~\ref{fig:seq_uc01_n1_operator}.

\begin{figure}[htbp]
\centering
\resizebox{\textwidth}{!}{%
\footnotesize
\begin{sequencediagram}
    \newthread{op}{op:OperatorTernak}
    \newinst[1.5]{ui}{ui:DashboardView}
    \newinst[1.5]{cs}{cs:Cage}
    \newinst[1.5]{tk}{tk:Task}
    \newinst[1.5]{db}{db:Database}

    \begin{call}{op}{Buka Halaman Dasbor Ternak}{ui}{Menampilkan ringkasan populasi, tugas, \& alarm kehamilan}
        \begin{call}{ui}{GetCageStats()}{cs}{Data statistik kandang}
            \begin{call}{cs}{GetCageStats()}{db}{Rows statistik}
            \end{call}
        \end{call}
        \begin{call}{ui}{GetMyTasks(id, role, date)}{tk}{Daftar tugas harian hari ini}
            \begin{call}{tk}{GetMyTasks(id, role, date)}{db}{Rows tugas}
            \end{call}
        \end{call}
    \end{call}

    \begin{call}{op}{Klik salah satu tugas harian pada daftar tugas}{ui}{Formulir pencatatan terkait ter-prefill}
        \begin{call}{ui}{GetMyTasks(id, role, date)}{tk}{Data detail tugas}
            \begin{call}{tk}{GetMyTasks(id, role, date)}{db}{Rows tugas}
            \end{call}
        \end{call}
    \end{call}
\end{sequencediagram}%
}
\caption{Sequence Diagram UC-01 N-1: Operator Melihat Dasbor Ternak}
\label{fig:seq_uc01_n1_operator}
\end{figure}

\newpage

% -------------------------------------------------------------------------
% UC-01 Normal-2: Admin melihat Dasbor Peternakan
% -------------------------------------------------------------------------
\item{Sequence Diagram UC-01 N-2: Admin Melihat Dasbor Peternakan}

Admin membuka halaman Dasbor Peternakan. Antarmuka memanggil
\texttt{GetCageStats()} pada kelas \texttt{cs:Cage} untuk menampilkan
ringkasan total populasi, jumlah kandang, dan persentase utilitas kapasitas.
Selanjutnya, sistem memanggil \texttt{GetSubmissions("pending")} pada kelas
\texttt{sb:Submission} untuk memuat seluruh pengajuan pencatatan Operator
yang belum mendapat persetujuan. Kedua data tersebut ditampilkan bersama
dengan lencana jumlah pengajuan dan alat ekspor laporan. Apabila Admin
mengeklik tombol Ekspor Laporan, sistem mengunduh laporan aktivitas
peternakan dalam format dokumen eksternal. Alur ini disajikan pada
Gambar~\ref{fig:seq_uc01_n2_admin}.

\begin{figure}[htbp]
\centering
\resizebox{\textwidth}{!}{%
\footnotesize
\begin{sequencediagram}
    \newthread{adm}{adm:Admin}
    \newinst[1.5]{ui}{ui:DashboardView}
    \newinst[1.5]{cs}{cs:Cage}
    \newinst[1.5]{sb}{sb:Submission}
    \newinst[1.5]{db}{db:Database}

    \begin{call}{adm}{Buka Halaman Dasbor Peternakan}{ui}{Menampilkan statistik, pengajuan pending, \& alat ekspor}
        \begin{call}{ui}{GetCageStats()}{cs}{Statistik total populasi \& utilitas kandang}
            \begin{call}{cs}{GetCageStats()}{db}{Rows statistik}
            \end{call}
        \end{call}
        \begin{call}{ui}{GetSubmissions("pending")}{sb}{Daftar pengajuan menunggu persetujuan}
            \begin{call}{sb}{GetSubmissions("pending")}{db}{Rows submissions berstatus pending}
            \end{call}
        \end{call}
    \end{call}

    \begin{call}{adm}{Klik tombol Ekspor Laporan}{ui}{File laporan terunduh}
    \end{call}
\end{sequencediagram}%
}
\caption{Sequence Diagram UC-01 N-2: Admin Melihat Dasbor Peternakan}
\label{fig:seq_uc01_n2_admin}
\end{figure}

\newpage

% -------------------------------------------------------------------------
% UC-01 Normal-3: Pemilik melihat Dasbor Pemilik (read-only)
% -------------------------------------------------------------------------
\item{Sequence Diagram UC-01 N-3: Pemilik Melihat Dasbor (Read-Only)}

Pemilik membuka halaman Dasbor yang menyajikan data eksekutif peternakan
bersifat \textit{read-only}. Sistem memanggil \texttt{GetCageStats()} pada
kelas \texttt{cs:Cage} untuk menampilkan metrik performa ringkasan seperti
rata-rata ADG, jumlah mutasi, dan total populasi. Sistem juga memanggil
\texttt{GetCageWeightStats()} untuk memuat data bobot kandang, tren produksi
pupuk kandang, dan tabel utilitas kapasitas kandang. Halaman dasbor Pemilik
tidak menyertakan tombol aksi pencatatan maupun tombol ekspor laporan,
sehingga seluruh tampilan bersifat hanya baca. Alur ini disajikan pada
Gambar~\ref{fig:seq_uc01_n3_pemilik}.

\begin{figure}[htbp]
\centering
\resizebox{\textwidth}{!}{%
\footnotesize
\begin{sequencediagram}
    \newthread{own}{own:Pemilik}
    \newinst[2]{ui}{ui:DashboardView}
    \newinst[2]{cs}{cs:Cage}
    \newinst[2]{db}{db:Database}

    \begin{call}{own}{Buka Halaman Dasbor Pemilik}{ui}{Menampilkan data eksekutif peternakan (read-only)}
        \begin{call}{ui}{GetCageStats()}{cs}{Metrik performa: ADG, mutasi, populasi}
            \begin{call}{cs}{GetCageStats()}{db}{Rows statistik}
            \end{call}
        \end{call}
        \begin{call}{ui}{GetCageWeightStats()}{cs}{Data bobot, tren produksi, utilitas kapasitas}
            \begin{call}{cs}{GetCageWeightStats()}{db}{Rows statistik bobot kandang}
            \end{call}
        \end{call}
    \end{call}
\end{sequencediagram}%
}
\caption{Sequence Diagram UC-01 N-3: Pemilik Melihat Dasbor (Read-Only)}
\label{fig:seq_uc01_n3_pemilik}
\end{figure}

\newpage

% -------------------------------------------------------------------------
% UC-01 Tidak Normal-1: Data dasbor gagal dimuat
% -------------------------------------------------------------------------
\item{Sequence Diagram UC-01 TN-1: Data Dasbor Gagal Dimuat}

Pengguna membuka halaman Dasbor, namun pemanggilan \texttt{GetCageStats()}
pada kelas \texttt{cs:Cage} gagal karena gangguan koneksi jaringan atau
server \textit{backend} tidak merespons. Kelas \texttt{cs:Cage} meneruskan
permintaan ke \texttt{db:Database}, tetapi basis data mengembalikan pesan
\textit{connection timeout}. Kesalahan ini dipropagasikan kembali ke
antarmuka, yang kemudian menampilkan pesan peringatan kepada pengguna
bahwa data gagal dimuat dari \textit{backend}. Alur ini disajikan pada
Gambar~\ref{fig:seq_uc01_tn1_error}.

\begin{figure}[htbp]
\centering
\resizebox{\textwidth}{!}{%
\footnotesize
\begin{sequencediagram}
    \newthread{usr}{usr:Pengguna}
    \newinst[2]{ui}{ui:DashboardView}
    \newinst[2]{cs}{cs:Cage}
    \newinst[2]{db}{db:Database}

    \begin{call}{usr}{Buka Halaman Dasbor}{ui}{Tampilkan pesan ``Gagal memuat data dari backend''}
        \begin{call}{ui}{GetCageStats()}{cs}{Error: koneksi gagal}
            \begin{call}{cs}{GetCageStats()}{db}{Connection timeout}
            \end{call}
        \end{call}
    \end{call}
\end{sequencediagram}%
}
\caption{Sequence Diagram UC-01 TN-1: Data Dasbor Gagal Dimuat}
\label{fig:seq_uc01_tn1_error}
\end{figure}

\newpage

% =========================================================================
% UC-02: KELOLA DATA PERKAWINAN
% =========================================================================

% -------------------------------------------------------------------------
% UC-02 Normal-1 & TN-1 & TN-2: Perkawinan + sdblock alt CoI
% -------------------------------------------------------------------------
\item{Sequence Diagram UC-02: Kelola Data Perkawinan}

Operator membuka menu Pencatatan Perkawinan. Sistem memanggil
\texttt{GetSheepList(filter)} pada kelas \texttt{ss:Sheep} untuk mengisi
dropdown daftar domba betina berstatus birahi dan pejantan aktif. Setelah
Operator memilih ID domba betina dan ID pejantan, sistem secara
\textit{real-time} memanggil \texttt{GetSheepGenealogy()} untuk menelusuri
silsilah leluhur kedua domba hingga 5 generasi ke atas guna menghitung
Koefisien Inbreeding (CoI) menggunakan metode jalur Wright.

Diagram menampilkan dua skenario menggunakan blok \textit{alt}. Jika CoI
$\ge$ 6.25\%, antarmuka menampilkan tanda bahaya merah dan memblokir aksi
penyimpanan, sehingga Operator tidak dapat melanjutkan. Jika CoI~$<$~6.25\%
atau silsilah tidak lengkap (dengan peringatan kuning), Operator mengisi
tanggal kawin dan catatan kemudian menekan Simpan. Sistem memanggil
\texttt{RecordMating(m)} pada kelas \texttt{ms:Mating} untuk menyimpan
data perkawinan, dilanjutkan dengan memanggil \texttt{CreateSubmission(sub)}
pada kelas \texttt{sb:Submission} untuk mendaftarkan pengajuan berstatus
\textit{pending} yang menunggu persetujuan Admin. Alur ini disajikan pada
Gambar~\ref{fig:seq_uc02_perkawinan}.

\begin{figure}[htbp]
\centering
\resizebox{\textwidth}{!}{%
\footnotesize
\begin{sequencediagram}
    \newthread{op}{op:OperatorTernak}
    \newinst[1.2]{ui}{ui:RecordFormView}
    \newinst[1.2]{ss}{ss:Sheep}
    \newinst[1.2]{ms}{ms:Mating}
    \newinst[1.2]{sb}{sb:Submission}
    \newinst[1.2]{db}{db:Database}

    % -- Muat daftar domba --
    \begin{call}{op}{Buka menu Pencatatan Perkawinan}{ui}{Tampilkan form \& dropdown domba}
        \begin{call}{ui}{GetSheepList(filter)}{ss}{Daftar domba betina birahi \& pejantan aktif}
            \begin{call}{ss}{GetSheepList(filter)}{db}{Rows domba}
            \end{call}
        \end{call}
    \end{call}

    % -- Cek silsilah secara real-time --
    \begin{call}{op}{Pilih ID domba betina \& ID pejantan}{ui}{Tampilkan status risiko inbreeding}
        \begin{call}{ui}{GetSheepGenealogy()}{ss}{Nilai CoI \& status kekerabatan}
            \begin{call}{ss}{GetSheepGenealogy()}{db}{Data leluhur 5 generasi}
            \end{call}
        \end{call}
    \end{call}

    % -- Percabangan berdasarkan hasil CoI --
    \begin{sdblock}{alt}{CoI $\ge$ 6.25\% -- Inbreeding Tinggi}
        \begin{call}{op}{Klik tombol ``Simpan''}{ui}{Blokir aksi -- Tampilkan pop-up error}
        \end{call}
        \mess{ui}{Tampilkan pesan: ``Risiko Inbreeding Tinggi: Pencatatan Diblokir''}{op}
    \end{sdblock}

    \begin{sdblock}{alt}{CoI < 6.25\% atau Silsilah Tidak Lengkap (Peringatan)}
        \begin{call}{op}{Isi tanggal kawin \& catatan, klik ``Simpan''}{ui}{Pesan ``Data Perkawinan Terkirim -- Menunggu Persetujuan''}
            \begin{call}{ui}{RecordMating(m)}{ms}{Data perkawinan tersimpan}
                \begin{call}{ms}{RecordMating(m)}{db}{Insert ke tabel matings}
                \end{call}
            \end{call}
            \begin{call}{ui}{CreateSubmission(sub)}{sb}{Pengajuan dibuat berstatus pending}
                \begin{call}{sb}{CreateSubmission(sub)}{db}{Insert ke tabel submissions}
                \end{call}
            \end{call}
        \end{call}
    \end{sdblock}
\end{sequencediagram}%
}
\caption{Sequence Diagram UC-02: Kelola Data Perkawinan}
\label{fig:seq_uc02_perkawinan}
\end{figure}

\newpage

% -------------------------------------------------------------------------
% UC-02 Tidak Normal-3: Formulir tidak lengkap (validasi gagal)
% -------------------------------------------------------------------------
\item{Sequence Diagram UC-02 TN-3: Formulir Perkawinan Tidak Lengkap}

Operator mengosongkan salah satu kolom wajib pada formulir pencatatan
perkawinan, seperti tanggal kawin atau nomor batch semen untuk
Inseminasi Buatan (IB), kemudian menekan tombol Simpan. Antarmuka
memanggil \texttt{RecordMating(m)} pada kelas \texttt{ms:Mating}, namun
kelas tersebut mengembalikan kesalahan validasi karena data masukan
tidak lengkap. Sistem menampilkan pesan kesalahan spesifik kepada
Operator, seperti ``Tanggal Kawin wajib diisi'' atau ``Asal Semen wajib
diisi'', dan membatalkan pengiriman formulir. Alur ini disajikan pada
Gambar~\ref{fig:seq_uc02_tn3_validasi}.

\begin{figure}[htbp]
\centering
\resizebox{\textwidth}{!}{%
\footnotesize
\begin{sequencediagram}
    \newthread{op}{op:OperatorTernak}
    \newinst[2.5]{ui}{ui:RecordFormView}
    \newinst[2.5]{ms}{ms:Mating}

    \begin{call}{op}{Kosongkan tanggal kawin / nomor batch semen}{ui}{}
    \end{call}

    \begin{call}{op}{Klik tombol ``Simpan''}{ui}{Tampilkan pesan validasi -- pengiriman dibatalkan}
        \begin{call}{ui}{RecordMating(m)}{ms}{Error: validasi masukan gagal}
        \end{call}
    \end{call}
    \mess{ui}{Tampilkan pesan: ``Tanggal Kawin wajib diisi'' / ``Asal Semen wajib diisi''}{op}
\end{sequencediagram}%
}
\caption{Sequence Diagram UC-02 TN-3: Formulir Perkawinan Tidak Lengkap}
\label{fig:seq_uc02_tn3_validasi}
\end{figure}

\newpage

% =========================================================================
% UC-03: KELOLA DATA BIRAHI
% =========================================================================

% -------------------------------------------------------------------------
% UC-03 Normal-1 & Normal-2: Pencatatan birahi + sdblock alt hasil
% -------------------------------------------------------------------------
\item{Sequence Diagram UC-03: Kelola Data Birahi}

Operator membuka menu Pencatatan Birahi. Sistem memanggil
\texttt{GetSheepList(filter)} pada kelas \texttt{ss:Sheep} untuk mengisi
dropdown pilihan domba betina yang tersedia. Operator memilih domba
betina, mengisi tanggal pemeriksaan, lalu memilih hasil cek birahi.

Diagram menampilkan dua skenario menggunakan blok \textit{alt}. Jika
Operator memilih hasil ``Birahi (Siap Kawin)'', sistem memanggil
\texttt{RecordEstrusCheck(ec)} pada kelas \texttt{ms:Mating} dan
\texttt{CreateSubmission(sub)} pada kelas \texttt{sb:Submission}. Setelah
disetujui Admin, status kesiapan kawin domba tersebut aktif dan namanya
muncul di \textit{dropdown} calon induk pada modul perkawinan. Jika
Operator memilih ``Tidak Birahi'', alur pencatatan dan pengajuan identik,
namun setelah disetujui, domba tersebut dihapus dari daftar acuan siap
kawin. Alur ini disajikan pada Gambar~\ref{fig:seq_uc03_birahi}.

\begin{figure}[htbp]
\centering
\resizebox{\textwidth}{!}{%
\footnotesize
\begin{sequencediagram}
    \newthread{op}{op:OperatorTernak}
    \newinst[1.2]{ui}{ui:RecordFormView}
    \newinst[1.2]{ss}{ss:Sheep}
    \newinst[1.2]{ms}{ms:Mating}
    \newinst[1.2]{sb}{sb:Submission}
    \newinst[1.2]{db}{db:Database}

    % -- Muat daftar domba betina --
    \begin{call}{op}{Buka menu Pencatatan Birahi}{ui}{Tampilkan form \& dropdown domba betina}
        \begin{call}{ui}{GetSheepList(filter)}{ss}{Daftar domba betina}
            \begin{call}{ss}{GetSheepList(filter)}{db}{Rows domba betina}
            \end{call}
        \end{call}
    \end{call}

    % -- Percabangan berdasarkan hasil cek birahi --
    \begin{sdblock}{alt}{Hasil: ``Birahi (Siap Kawin)''}
        \begin{call}{op}{Pilih domba, isi tanggal, pilih ``Birahi'', klik Simpan}{ui}{Pesan ``Catatan Birahi Berhasil Diajukan''}
            \begin{call}{ui}{RecordEstrusCheck(ec)}{ms}{Data cek birahi tersimpan}
                \begin{call}{ms}{RecordEstrusCheck(ec)}{db}{Insert ke tabel estrus\_checks}
                \end{call}
            \end{call}
            \begin{call}{ui}{CreateSubmission(sub)}{sb}{Pengajuan dibuat berstatus pending}
                \begin{call}{sb}{CreateSubmission(sub)}{db}{Insert ke tabel submissions}
                \end{call}
            \end{call}
        \end{call}
        \mess{ui}{Domba aktif muncul di daftar calon induk modul perkawinan}{op}
    \end{sdblock}

    \begin{sdblock}{alt}{Hasil: ``Tidak Birahi''}
        \begin{call}{op}{Pilih domba, isi tanggal, pilih ``Tidak Birahi'', klik Simpan}{ui}{Pesan ``Catatan Birahi Berhasil Diajukan''}
            \begin{call}{ui}{RecordEstrusCheck(ec)}{ms}{Data cek birahi tersimpan}
                \begin{call}{ms}{RecordEstrusCheck(ec)}{db}{Insert ke tabel estrus\_checks}
                \end{call}
            \end{call}
            \begin{call}{ui}{CreateSubmission(sub)}{sb}{Pengajuan dibuat berstatus pending}
                \begin{call}{sb}{CreateSubmission(sub)}{db}{Insert ke tabel submissions}
                \end{call}
            \end{call}
        \end{call}
        \mess{ui}{Domba dihapus dari daftar acuan siap kawin di modul perkawinan}{op}
    \end{sdblock}
\end{sequencediagram}%
}
\caption{Sequence Diagram UC-03: Kelola Data Birahi}
\label{fig:seq_uc03_birahi}
\end{figure}

\newpage

% -------------------------------------------------------------------------
% UC-03 Tidak Normal-1: Data tidak lengkap saat menyimpan
% -------------------------------------------------------------------------
\item{Sequence Diagram UC-03 TN-1: Data Birahi Tidak Lengkap}

Operator menekan tombol Simpan tanpa memilih ID domba betina atau tanpa
menentukan hasil cek birahi. Antarmuka memanggil
\texttt{RecordEstrusCheck(ec)} pada kelas \texttt{ms:Mating}, namun
kelas tersebut mengembalikan kesalahan validasi karena data masukan
tidak lengkap. Sistem menampilkan pesan kesalahan spesifik kepada
Operator, yaitu ``ID Domba wajib dipilih'' atau ``Hasil Cek Birahi wajib
dipilih'', dan membatalkan pengiriman formulir. Alur ini disajikan pada
Gambar~\ref{fig:seq_uc03_tn1_validasi}.

\begin{figure}[htbp]
\centering
\resizebox{\textwidth}{!}{%
\footnotesize
\begin{sequencediagram}
    \newthread{op}{op:OperatorTernak}
    \newinst[2.5]{ui}{ui:RecordFormView}
    \newinst[2.5]{ms}{ms:Mating}

    \begin{call}{op}{Klik ``Simpan'' tanpa memilih ID domba / hasil cek birahi}{ui}{Tampilkan pesan validasi -- pengiriman dibatalkan}
        \begin{call}{ui}{RecordEstrusCheck(ec)}{ms}{Error: validasi masukan gagal}
        \end{call}
    \end{call}
    \mess{ui}{Tampilkan pesan: ``ID Domba wajib dipilih'' / ``Hasil Cek Birahi wajib dipilih''}{op}
\end{sequencediagram}%
}
\caption{Sequence Diagram UC-03 TN-1: Data Birahi Tidak Lengkap}
\label{fig:seq_uc03_tn1_validasi}
\end{figure}

\newpage

% =========================================================================
% UC-04: KELOLA DATA KONVERSI PAKAN
% =========================================================================

% -------------------------------------------------------------------------
% UC-04 Normal-1: Inisiasi konversi silase baru (dengan sdblock loop)
% -------------------------------------------------------------------------
\item{Sequence Diagram UC-04 N-1: Inisiasi Konversi Pakan Silase Baru}

Operator membuka formulir Konversi Pakan. Sistem memanggil
\texttt{GetMasterFeedList()} pada kelas \texttt{fs:Feed} untuk menampilkan
daftar bahan baku pakan mentah beserta stok yang tersedia, dikelompokkan
per kategori nutrisi (Serat/Hijauan, Energi, Protein, dan
Aktivator/Mineral). Operator mengisi tanggal konversi, target kuantitas
silase, dan memilih bahan baku dari masing-masing kategori.

Setelah Operator menekan Simpan, antarmuka memanggil
\texttt{RecordSilageConversion(sc)} pada kelas \texttt{fs:Feed}. Di dalam
proses ini, sistem mengiterasi setiap bahan baku yang digunakan
menggunakan blok \textit{loop}, kemudian memanggil
\texttt{UpdateFeedStock(amount, "kurang")} ke basis data untuk memotong
stok masing-masing bahan baku mentah. Setelah itu, kelas \texttt{fs:Feed}
menyimpan catatan konversi baru dengan status awal ``fermentasi'' dan
target siklus 21 hari. Terakhir, antarmuka memanggil \texttt{CreateTask(t)}
pada kelas \texttt{tk:Task} untuk mendaftarkan tugas pengecekan fermentasi
silase secara otomatis 7 hari dari tanggal pembuatan. Alur ini disajikan
pada Gambar~\ref{fig:seq_uc04_n1_konversi}.

\begin{figure}[htbp]
\centering
\resizebox{\textwidth}{!}{%
\footnotesize
\begin{sequencediagram}
    \newthread{op}{op:OperatorTernak}
    \newinst[1.5]{ui}{ui:RecordFormView}
    \newinst[1.5]{fs}{fs:Feed}
    \newinst[1.5]{tk}{tk:Task}
    \newinst[1.5]{db}{db:Database}

    % -- Muat daftar bahan baku pakan --
    \begin{call}{op}{Buka form Konversi Pakan}{ui}{Tampilkan form \& dropdown bahan baku per kategori}
        \begin{call}{ui}{GetMasterFeedList()}{fs}{Daftar bahan baku pakan beserta stok}
            \begin{call}{fs}{GetMasterFeedList()}{db}{Rows master pakan}
            \end{call}
        \end{call}
    \end{call}

    % -- Simpan konversi silase baru --
    \begin{call}{op}{Isi tanggal, target qty, pilih bahan baku 4 kategori, klik Simpan}{ui}{Pesan ``Proses Konversi Berhasil Dimulai''}
        \begin{call}{ui}{RecordSilageConversion(sc)}{fs}{Record fermentasi terinisiasi}
            \begin{call}{fs}{RecordSilageConversion(sc)}{db}{Insert record konversi berstatus ``fermentasi''}
            \end{call}
        \end{call}
        \begin{sdblock}{loop}{untuk setiap bahan baku (Serat, Energi, Protein, Aktivator)}
            \begin{call}{ui}{UpdateFeedStock(amount, "kurang")}{fs}{Stok terupdate}
                \begin{call}{fs}{UpdateFeedStock(amount, "kurang")}{db}{Kurangi stok bahan baku mentah}
                \end{call}
            \end{call}
        \end{sdblock}
        \begin{call}{ui}{CreateTask(t)}{tk}{Tugas pengecekan fermentasi terjadwal}
            \begin{call}{tk}{CreateTask(t)}{db}{Insert tugas (tgl\_sekarang + 7 hari)}
            \end{call}
        \end{call}
    \end{call}
\end{sequencediagram}%
}
\caption{Sequence Diagram UC-04 N-1: Inisiasi Konversi Pakan Silase Baru}
\label{fig:seq_uc04_n1_konversi}
\end{figure}

\newpage

% -------------------------------------------------------------------------
% UC-04 Normal-2: Operator memantau daftar fermentasi silase
% -------------------------------------------------------------------------
\item{Sequence Diagram UC-04 N-2: Memantau Daftar Fermentasi Silase}

Operator membuka menu Gudang dan berpindah ke tab Fermentasi untuk
memantau seluruh riwayat konversi pakan silase. Sistem memanggil
\texttt{GetFermentationLogs(id)} pada kelas \texttt{fer:SilageFermentationLog}
yang meneruskan permintaan ke \texttt{db:Database} untuk mengambil
seluruh data konversi beserta log perkembangan yang telah dicatat.
Antarmuka menampilkan daftar konversi yang memuat persentase
\textit{progress} fermentasi (dihitung berdasarkan hari ke-X dari target
21 hari), komposisi bahan baku yang digunakan, dan status terkini setiap
batch. Alur ini disajikan pada Gambar~\ref{fig:seq_uc04_n2_pantau}.

\begin{figure}[htbp]
\centering
\resizebox{\textwidth}{!}{%
\footnotesize
\begin{sequencediagram}
    \newthread{op}{op:OperatorTernak}
    \newinst[2]{ui}{ui:WarehouseView}
    \newinst[2]{fer}{fer:SilageFermentationLog}
    \newinst[2]{db}{db:Database}

    \begin{call}{op}{Buka menu Gudang $\rightarrow$ tab Fermentasi}{ui}{Tampilkan daftar: progress \%, komposisi, status}
        \begin{call}{ui}{GetFermentationLogs(id)}{fer}{Seluruh data konversi \& log perkembangan}
            \begin{call}{fer}{GetFermentationLogs(id)}{db}{Rows konversi \& log fermentasi}
            \end{call}
        \end{call}
    \end{call}
\end{sequencediagram}%
}
\caption{Sequence Diagram UC-04 N-2: Operator Memantau Daftar Fermentasi Silase}
\label{fig:seq_uc04_n2_pantau}
\end{figure}

\newpage

% -------------------------------------------------------------------------
% UC-04 Tidak Normal-1: Komposisi bahan baku tidak lengkap
% -------------------------------------------------------------------------
\item{Sequence Diagram UC-04 TN-1: Komposisi Bahan Baku Tidak Lengkap}

Operator mengosongkan salah satu dari empat kategori bahan baku yang
wajib diisi, misalnya kategori Aktivator/Mineral dikosongkan, kemudian
menekan tombol Simpan. Antarmuka memanggil \texttt{RecordSilageConversion(sc)}
pada kelas \texttt{fs:Feed}, namun kelas tersebut mengembalikan kesalahan
karena komposisi bahan baku tidak memenuhi syarat minimum empat kelompok
nutrisi. Sistem memblokir penyimpanan dan menampilkan pesan kesalahan
kepada Operator: ``Semua kelompok bahan baku (Serat, Energi, Protein,
Aktivator) wajib diisi''. Alur ini disajikan pada
Gambar~\ref{fig:seq_uc04_tn1_validasi}.

\begin{figure}[htbp]
\centering
\resizebox{\textwidth}{!}{%
\footnotesize
\begin{sequencediagram}
    \newthread{op}{op:OperatorTernak}
    \newinst[2.5]{ui}{ui:RecordFormView}
    \newinst[2.5]{fs}{fs:Feed}

    \begin{call}{op}{Kosongkan salah satu kategori bahan baku (misal: Aktivator)}{ui}{}
    \end{call}

    \begin{call}{op}{Klik tombol Simpan}{ui}{Tampilkan pesan validasi -- penyimpanan diblokir}
        \begin{call}{ui}{RecordSilageConversion(sc)}{fs}{Error: komposisi bahan baku tidak lengkap}
        \end{call}
    \end{call}
    \mess{ui}{Tampilkan pesan: ``Semua kelompok bahan baku (Serat, Energi, Protein, Aktivator) wajib diisi''}{op}
\end{sequencediagram}%
}
\caption{Sequence Diagram UC-04 TN-1: Komposisi Bahan Baku Konversi Pakan Tidak Lengkap}
\label{fig:seq_uc04_tn1_validasi}
\end{figure}

\newpage

% =========================================================================
% UC-05: KELOLA DATA STATUS KONVERSI PAKAN
% =========================================================================

% -------------------------------------------------------------------------
% UC-05 Normal-1/2/3: Log fermentasi + sdblock alt status
% -------------------------------------------------------------------------
\item{Sequence Diagram UC-05: Kelola Data Status Konversi Pakan}

Operator membuka menu Gudang pada tab Fermentasi dan mengeklik tombol
``Update Progres'' pada batch konversi silase yang dituju. Sistem
memanggil \texttt{GetFermentationLogs(id)} pada kelas
\texttt{fer:SilageFermentationLog} untuk mengambil seluruh riwayat log
perkembangan fermentasi yang telah tercatat dan menampilkannya di dalam
modal. Operator kemudian mengisi nilai pH, suhu wadah, catatan kondisi
fisik, dan memilih status kesiapan fermentasi.

Setelah Operator menekan Simpan, sistem memanggil
\texttt{CreateFermentationLog(l)} pada kelas
\texttt{fer:SilageFermentationLog} untuk menyimpan log pengecekan baru ke
basis data. Selanjutnya, diagram menampilkan percabangan menggunakan blok
\textit{alt}: jika status yang dipilih adalah ``Selesai (Siap Disajikan)'',
sistem memanggil \texttt{UpdateFeedStock(amount, "tambah")} pada kelas
\texttt{fs:Feed} untuk secara otomatis menambah kuantitas target pakan
silase ke stok gudang utama; jika status yang dipilih adalah ``Gagal
(Buang)'' atau ``Masih Proses Fermentasi'', tidak ada perubahan stok dan
status konversi diperbarui sesuai pilihan Operator. Alur ini disajikan
pada Gambar~\ref{fig:seq_uc05_status_konversi}.

\begin{figure}[htbp]
\centering
\resizebox{\textwidth}{!}{%
\footnotesize
\begin{sequencediagram}
    \newthread{op}{op:OperatorTernak}
    \newinst[1.2]{ui}{ui:WarehouseView}
    \newinst[1.2]{fer}{fer:SilageFermentationLog}
    \newinst[1.2]{fs}{fs:Feed}
    \newinst[1.2]{db}{db:Database}

    % -- Ambil log sebelumnya & tampilkan modal --
    \begin{call}{op}{Klik ``Update Progres'' pada konversi yang dituju}{ui}{Tampilkan modal Kelola Fermentasi}
        \begin{call}{ui}{GetFermentationLogs(id)}{fer}{Daftar log perkembangan sebelumnya}
            \begin{call}{fer}{GetFermentationLogs(id)}{db}{Rows log fermentasi}
            \end{call}
        \end{call}
    \end{call}

    % -- Simpan log baru --
    \begin{call}{op}{Input pH, suhu, catatan, pilih status \& klik ``Simpan''}{ui}{Pesan ``Status Fermentasi Berhasil Diperbarui''}
        \begin{call}{ui}{CreateFermentationLog(l)}{fer}{Log baru tersimpan}
            \begin{call}{fer}{CreateFermentationLog(l)}{db}{Insert log baru ke DB}
            \end{call}
        \end{call}
    \end{call}

    % -- Percabangan berdasarkan Status Kesiapan --
    \begin{sdblock}{opt}{Status: ``Selesai (Siap Disajikan)''}
        \begin{call}{ui}{UpdateFeedStock(amount, "tambah")}{fs}{Stok pakan silase bertambah di gudang}
            \begin{call}{fs}{UpdateFeedStock(amount, "tambah")}{db}{Update stok pakan silase}
            \end{call}
        \end{call}
    \end{sdblock}
\end{sequencediagram}%
}
\caption{Sequence Diagram UC-05: Kelola Data Status Konversi Pakan}
\label{fig:seq_uc05_status_konversi}
\end{figure}

\newpage

% -------------------------------------------------------------------------
% UC-05 Tidak Normal-1: Radio status tidak dipilih
% -------------------------------------------------------------------------
\item{Sequence Diagram UC-05 TN-1: Status Kesiapan Tidak Dipilih}

Operator memasukkan nilai pH dan suhu wadah pada modal pengecekan
fermentasi, tetapi tidak memilih pilihan radio status kesiapan fermentasi,
lalu menekan tombol Simpan. Antarmuka memanggil
\texttt{CreateFermentationLog(l)} pada kelas
\texttt{fer:SilageFermentationLog}, namun kelas tersebut mengembalikan
kesalahan validasi karena status kesiapan merupakan kolom wajib yang
harus diisi. Sistem memblokir penyimpanan dan menampilkan pesan kesalahan
kepada Operator: ``Status kesiapan fermentasi wajib dipilih''. Alur ini
disajikan pada Gambar~\ref{fig:seq_uc05_tn1_validasi}.

\begin{figure}[htbp]
\centering
\resizebox{\textwidth}{!}{%
\footnotesize
\begin{sequencediagram}
    \newthread{op}{op:OperatorTernak}
    \newinst[2.5]{ui}{ui:WarehouseView}
    \newinst[2.5]{fer}{fer:SilageFermentationLog}

    \begin{call}{op}{Input pH \& suhu -- tidak memilih radio status kesiapan}{ui}{}
    \end{call}

    \begin{call}{op}{Klik tombol Simpan}{ui}{Tampilkan pesan error -- penyimpanan diblokir}
        \begin{call}{ui}{CreateFermentationLog(l)}{fer}{Error: status kesiapan wajib dipilih}
        \end{call}
    \end{call}
    \mess{ui}{Tampilkan pesan error: ``Status kesiapan fermentasi wajib dipilih''}{op}
\end{sequencediagram}%
}
\caption{Sequence Diagram UC-05 TN-1: Status Kesiapan Fermentasi Tidak Dipilih}
\label{fig:seq_uc05_tn1_validasi}
\end{figure}

\end{document}
